% Slides 126--145
\sectionframe{07 / 7--14 April 2026}{The quantitative sandbox becomes an operated web application}{Commits 7e08907, 65a2acb, 711ce0d, b9374c8, 754707f, 297e02c, d3b45a4.}

\lesson{The application factory separates creation from execution}
{\texttt{create\_app()} loads configuration, registers routes, and returns a Flask application.}
{\texttt{app.py} is the local entrypoint; \texttt{wsgi.py} is the production entrypoint.}
{Tests can create isolated apps with temporary configuration.}
{This pattern prevents server startup side effects from contaminating model tests.}
{tgir\_quantlib\_tools/app\_factory.py; app.py; wsgi.py.}

\lesson{Authentication protects the workstation boundary}
{The root page is a login screen and quantitative routes require a session.}
{Credentials and the Flask secret come from environment configuration.}
{Tests use explicit temporary credentials and disable secure cookies locally.}
{The repository does not commit production secrets.}
{tgir\_quantlib\_tools/auth.py and config.py; tests/test\_smoke.py.}

\lesson{Routes remain delivery adapters}
{Routes translate HTTP form or JSON input into normalized state.}
{Pricing remains in \texttt{portfolio.py} or the standalone XCCY module.}
{Context builders turn quantitative outputs into view-specific structures.}
{Keeping Flask thin lets CLI, tests, MCP, and future C++ clients reuse the same economics.}
{tgir\_quantlib\_tools/routes.py; dashboard.py.}

\lesson{Server-rendered templates fit the project's scale}
{Jinja templates produce the dashboard, trade forms, model page, and later XCCY lab.}
{There is no separate React build, API gateway, or database.}
{The architecture favors a single host and minimal infrastructure.}
{This is a conscious deviation, not an unfinished split frontend.}
{AGENTS.md; ARCHITECTURE.md.}

\lesson{Session state supports a demonstration workflow}
{User-edited market and trade inputs live in the Flask session.}
{Reset returns to deterministic defaults loaded from JSON.}
{A realtime tick perturbs the demo state and returns updated analytics.}
{Production multi-user state would require stronger persistence and audit semantics.}
{dashboard.py; routes.py; tests/test\_smoke.py.}

\chronology{11 Apr 2026}{UI and deployment documentation are consolidated}
{Commit \texttt{65a2acb} changes 31 files, adds a handoff, deployment guide, screenshots, and extensive UI refinements.}
{Local timestamps show the long-lived pricing module continuing to evolve after this commit.}
{Operational documentation becomes part of the delivered system.}
{A model is usable only when users can install, run, diagnose, and recover it.}
{Commit 65a2acb; DEPLOYMENT.md; SESSION\_HANDOFF.md.}

\lesson{CI encodes repeatable confidence}
{GitHub Actions creates Python 3.13, installs requirements, and runs selected test modules.}
{GitHub records successful CI runs for the April 2026 main-branch commits.}
{The local full suite is broader than the historical release-gate subset.}
{A green pipeline is meaningful only when its selected tests match current critical paths.}
{.github/workflows/ci.yml; GitHub Actions history.}

\lesson{DigitalOcean deployment is a simple three-tier stack}
{Apache terminates TLS and proxies requests.}
{Gunicorn binds the Flask app only on \texttt{127.0.0.1:8008}.}
{systemd owns process startup and restart under \texttt{www-data}; the application root is \texttt{/var/www/html/tgir\_quantlib\_tools}.}
{One approximately USD 7-per-month Ubuntu droplet keeps the hobby affordable while preserving explicit process boundaries.}
{DEPLOYMENT.md; deploy templates.}

\begin{frame}{Production request path}
  \centering
  \begin{tikzpicture}[node distance=9mm, every node/.style={font=\small}]
    \node[draw,rounded corners,fill=TgirMint!50,minimum width=25mm,minimum height=10mm] (browser) {Browser};
    \node[draw,rounded corners,fill=white,right=of browser,minimum width=29mm,minimum height=10mm] (apache) {Apache + TLS};
    \node[draw,rounded corners,fill=TgirMint!50,right=of apache,minimum width=27mm,minimum height=10mm] (gunicorn) {Gunicorn};
    \node[draw,rounded corners,fill=white,right=of gunicorn,minimum width=31mm,minimum height=10mm] (flask) {Flask + QuantLib};
    \draw[-{Latex},thick,TgirTeal] (browser) -- node[above]{HTTPS} (apache);
    \draw[-{Latex},thick,TgirTeal] (apache) -- node[above]{HTTP} (gunicorn);
    \draw[-{Latex},thick,TgirTeal] (gunicorn) -- (flask);
  \end{tikzpicture}
  \vspace{7mm}
  \begin{itemize}\setlength\itemsep{0.7em}
    \item Secrets remain in a mode-600 server-side \texttt{.env}; deployment builds a fresh \texttt{.venv} at its final path.
    \item Health checks test \texttt{127.0.0.1:8008/health} and \texttt{https://quant.tglauner.com/health}.
  \end{itemize}
  \takeaway{About USD 7 per month is enough to keep a serious quantitative hobby continuously reachable behind production-style controls.}
  \src{deploy/apache and deploy/systemd templates; DEPLOYMENT.md.}
\end{frame}

\lesson{Deployment failures expose architecture assumptions}
{GitHub shows successful CI but repeated failed DigitalOcean workflows in April 2026.}
{Automated delivery requires \texttt{DO\_HOST}, \texttt{DO\_USER}, \texttt{DO\_SSH\_KEY}, \texttt{DO\_APP\_ROOT}, and \texttt{DO\_HEALTHCHECK\_URL}.}
{A prebuilt virtualenv could not be moved because executable shebangs contained absolute paths; rollback succeeded.}
{The corrected release rebuilt its virtualenv at the final path and passed health checks.}
{GitHub Actions history; deployment performed 2026-08-21; README virtualenv note.}

\lesson{Rollback is part of deployment design}
{The live app is preserved before replacement.}
{A failed service health check restores the prior directory and restarts systemd.}
{The successful August deployment retained a timestamped rollback copy.}
{Recoverability matters more than pretending a release cannot fail.}
{DEPLOYMENT.md; August 2026 deployment verification.}

\chronology{13 Apr 2026}{An IBKR comparison dashboard is added}
{Commit \texttt{711ce0d} adds more than one thousand lines centered on a second dashboard.}
{The interface compares workbench values with an execution-platform-style view.}
{Routing and tests expand while the pricing core remains shared.}
{Multiple views can consume one normalized pricing state without duplicating model code.}
{Commit 711ce0d.}

\lesson{Presentation changes should not change economics}
{Subsequent commits tighten monitor layout, realtime ticks, and reset-control styling.}
{Commit names \texttt{754707f}, \texttt{297e02c}, and \texttt{d3b45a4} are primarily interface refinements.}
{Tests protect core marks and route behavior while CSS and templates evolve.}
{Chronology distinguishes visual iterations from model changes.}
{Git commits 754707f, 297e02c, d3b45a4.}

\lesson{A dashboard bug demonstrates signature drift}
{In August, \texttt{\_forward\_rate\_chart} required regular and scenario points but one caller supplied only one argument.}
{The dashboard returned HTTP 500 after login.}
{The call site was repaired and smoke tests were expanded.}
{Function signatures shared across UI paths need direct route tests.}
{Local traceback dated 2026-08-14; dashboard.py and tests/test\_smoke.py.}

\lesson{Local timestamps refine the August sequence}
{The XCCY quantitative summary was born 14 August.}
{The standalone pricer, JSON deal/market files, schemas, and core tests were born 15 August.}
{Diagnostics and MCP followed later on 15 August; validation and web work followed 16--17 August.}
{Because these landed in one Git commit, filesystem times provide supporting -- not authoritative -- ordering.}
{macOS birth/modification timestamps inspected 2026-08-21.}

\lesson{Configuration keeps paths and secrets outside pricing code}
{Config values identify market, deal, result, debug, authentication, and cookie behavior.}
{Tests override paths with temporary directories.}
{Production reads the same keys from the server environment.}
{The design avoids hard-coded personal filesystem paths in runtime logic.}
{tgir\_quantlib\_tools/config.py and tests.}

\lesson{Thread safety matters around QuantLib globals}
{The evaluation date and relinkable handles can be process-global or shared.}
{Concurrent valuations with different as-of dates can interfere if not isolated.}
{A production C++ or Python service should use serialized contexts or process isolation.}
{The current demo uses a small Gunicorn process model and is not a high-throughput risk grid.}
{REST integration specification; QuantLib Settings usage.}

\lesson{The architecture optimizes for inspectability}
{Source files remain direct and readable rather than distributed across services.}
{JSON snapshots, CSV exports, HTML diagnostics, CLI scripts, and tests expose intermediate states to Tim, Codex, and Claude alike.}
{The cost is limited concurrency and no database-backed audit store.}
{For a research workbench, this tradeoff enables two agents to build, review, extend validation, and reproduce each other's evidence.}
{AGENTS.md; ARCHITECTURE.md.}

\lesson{By mid-April, the platform is ready for a cross-asset challenge}
{Curves, rates options, equity options, calibrated trees, risk, tests, authentication, and deployment all exist.}
{The next requested product combines two curves, two short-rate factors, FX, and early cancellation.}
{The existing architecture offers reusable primitives but not a native engine.}
{Tim directs the August challenge; Codex begins with rigorous specifications before source implementation, creating a stable target for Claude's later review and validation extensions.}
{Git chronology through d3b45a4 and local August timestamps.}
